% linebreaking.tex — line breaking for languages that build long compounds.
%
% Usable on its own:  pandoc doc.md -H linebreaking.tex --pdf-engine=xelatex ...
%
% \tolerance raises how loose a line TeX will accept before it gives up and
% lets text run into the margin. The default is 200, which is strict enough
% that TeX will REFUSE an available hyphenation point if taking it would leave
% a loose line — so a break opportunity you inserted (see
% filters/linebreaks.lua) goes unused and the overfull line happens anyway.
%
% 1500 rather than TeX's default of 200. Measured on one corpus of German
% business documents, where 200 and 800 each left one overflowing line, 1500 left
% none, and 3000 changed nothing further — your corpus will differ, so treat the
% number as a considered default rather than a derived constant. The price is a
% handful of loose lines. Text in the margin is the worse fault.
\tolerance=1500

% NOTE ON \emergencystretch: pandoc's own LaTeX template already sets
% \setlength{\emergencystretch}{3em}. If you are using pandoc you have it
% whether you ask for it or not, and re-setting it here would be a no-op.
% Set it yourself only if you are not going through pandoc's template.
